\section{投资问题：从已知H1到未知H2}

盈利预测的核心不是给每个分部填满数字，而是用透明约束回答：H1已知高利润之后，H2和2027--2028年应保留多少。核心判断是，2026E归母77.5亿元已经承认文莱高景气，同时通过H2仅20.0亿元防止峰值年化；2027E回落至70.0亿元反映价差均值回归；2028E恢复至75.0亿元需要新项目和国内链改善，但不预先计入文莱二期完整贡献。

\begin{exhibitbox}[House盈利预测]
\noindent\begin{tabularx}{\textwidth}{L{3.4cm}R{2.0cm}R{2.0cm}R{2.0cm}R{2.0cm}}
\toprule
\textbf{亿元／元} & \textbf{2025A} & \textbf{2026E} & \textbf{2027E} & \textbf{2028E} \\
\midrule
营业收入 & 1,135.3 & 1,417.4 & 1,435.2 & 1,477.3 \\
House归母净利润 & 2.58 & 77.5 & 70.0 & 75.0 \\
公司披露会计EPS & 0.07\textsuperscript{*} & -- & -- & -- \\
估值稀释EPS & -- & 2.00 & 1.80 & 1.93 \\
国信归母净利润 & -- & 79.1 & 78.0 & 85.6 \\
\bottomrule
\end{tabularx}
\vspace{0.3em}
{\scriptsize \textsuperscript{*}2025A只作低谷对比；历史会计EPS受加权平均股本和实际转股时点影响，不能用期末完全稀释股本替代。}
\end{exhibitbox}
\sourcenote{HY-OFF-2025AR；HY-BRK-GUOXIN；analysis/earnings\_forecast.md；House预测截止2026-07-22。2026--2028收入采用国信原始PDF作规模检查，House归母为独立预测。}

2026E收入采用国信可审计原始报告作为总量检查，归母则由House独立构建。2026年77.5亿元低于国信79.1亿元，差异不大，但口径更保守：本报告使用38.80694189亿股已披露稀释分母，不回加未量化转债有效利息，并为员工计划股份支付未知费用保留缓冲。

\section{2025到2026：不构造伪精确分部瀑布}

2025A归母2.58亿元到2026E 77.5亿元，增量约74.9亿元。公开证据支持驱动顺序：文莱一期负荷与成品油裂解、PX／苯芳烃价差、国内PTA与聚酯加工差、广西CPL--PA6爬坡；抵消项包括利息、少数股东、库存与营运资金、可能检修和未量化员工费用。

由于H1预告没有项目净利润拆分，本报告不把74.9亿元强行分配到各板块。伪精确瀑布会让行业价差、公司利润和少数股东归属混在一起。当前最可靠的桥是总归母层面的H1--H2正常化，以及分部收入／毛利率对总量的约束。

\begin{exhibitbox}[盈利驱动与模型信用]
\small
\noindent\begin{tabularx}{\textwidth}{L{3.0cm}L{2.7cm}L{2.55cm}X}
\toprule
\textbf{驱动} & \textbf{2026证据} & \textbf{模型角色} & \textbf{不可做的推断} \\
\midrule
文莱成品油 & 满产满销、盈利高位 & 主要正向驱动 & 行业裂解乘全部销量 \\
PX／苯 & 公司称盈利高位 & 主要正向驱动 & 付费价差极值年化 \\
PTA／聚酯 & 预告称改善，行业修复不均 & 次要正向驱动 & 参控股产能乘加工费 \\
CPL／PA6 & 广西项目贡献改善 & 事件期权／缓冲 & 用设计产能推收入与利润 \\
少数股东 & Q1损益8.23亿元 & 归母折损 & 把合并利润等额归母 \\
利息／员工计划 & 有效利息与股份支付未量化 & 保守缓冲 & 未披露利润回加 \\
\bottomrule
\end{tabularx}
\end{exhibitbox}
\sourcenote{HY-OFF-2026H1P；HY-OFF-2026Q1；HY-OFF-ESOP7；analysis/chain\_earnings\_bridge.md，研究截止2026-07-22。}

\section{H1--H2桥：基准只保留20亿元}

H1预告中值57.5亿元，基准全年77.5亿元，意味着H2归母20.0亿元。熊／牛情景分别为全年65／95亿元，对应H2 7.5／37.5亿元。机械翻倍H1得到115亿元，被明确排除。H2假设综合反映产品价差、检修、原油库存、国内加工费、运费、汇率、少数股东和费用，而不是单一季节项。

\begin{exhibitbox}[H2敏感性与每股影响]
\noindent\begin{tabularx}{\textwidth}{L{2.4cm}R{2.0cm}R{2.0cm}R{2.0cm}X}
\toprule
\textbf{情景} & \textbf{全年归母} & \textbf{H2归母} & \textbf{稀释EPS} & \textbf{预测判断} \\
\midrule
熊 & 65.0亿元 & 7.5亿元 & 1.675元 & 高景气快速回落、现金和费用折损 \\
基准 & 77.5亿元 & 20.0亿元 & 1.997元 & 文莱正常化，国内温和修复 \\
牛 & 95.0亿元 & 37.5亿元 & 2.448元 & 文莱持续高景气、国内与现金共振 \\
机械年化 & 115.0亿元 & 57.5亿元 & 2.963元 & 拒绝：无周期、归属和费用校准 \\
\bottomrule
\end{tabularx}
\end{exhibitbox}
\sourcenote{HY-OFF-2026H1P；analysis/valuation\_audit.md。EPS按38.80694189亿股计算，未加员工计划现金；预测截止2026-07-22。}

H1预告55--60亿元扣除Q1已披露19.95亿元，隐含Q2归母35.05--40.05亿元，中值为37.55亿元。为消除“三季报时已经实现全年H2”的时间歧义，House进一步把H2拆为Q3E与Q4E；这只是节奏假设，不是公司指引，也不得用Q3机械外推全年。

\begin{exhibitbox}[Q2隐含值与Q3--Q4盈利桥]
\small
\noindent\begin{tabularx}{\textwidth}{L{2.2cm}R{1.7cm}R{1.7cm}R{1.7cm}R{1.8cm}X}
\toprule
\textbf{情景} & \textbf{Q2隐含} & \textbf{Q3E} & \textbf{Q4E} & \textbf{H2E} & \textbf{9M累计／解释} \\
\midrule
熊 & 35.05--40.05 & 4.5 & 3.0 & 7.5 & 62.0；价差快速回落 \\
基准 & 35.05--40.05 & 12.0 & 8.0 & 20.0 & 69.5；文莱正常化 \\
牛 & 35.05--40.05 & 22.5 & 15.0 & 37.5 & 80.0；高景气延续 \\
\bottomrule
\end{tabularx}
\end{exhibitbox}
\sourcenote{HY-OFF-2026Q1；HY-OFF-2026H1P；AStock House节奏假设，截止2026-07-22。Q2区间为H1预告减Q1；9M累计按H1中值57.5亿元加Q3E。}

每5亿元归母变化影响稀释EPS约0.129元，在7.5倍PE下约影响0.97元／股。该敏感性说明，15.7元目标对H2偏差并不迟钝；但真正的下行还可能来自利润下调与PE压缩同时发生。

\section{分部规模检查：收入合理，不代表利润确定}

为检验1,417.4亿元收入是否脱离现有业务结构，模型采用非发布型分部检查：文莱炼油约380亿元、芳烃化工140亿元、PTA 55亿元、PIA 15亿元、聚酯730亿元、尼龙35亿元、供应链服务62亿元，合计约1,417亿元。这不是公司指引，也不是正式分部预测，只用于识别总量错误。

在外部券商假设的炼油17.8\%、化工20\%、PTA 6\%、聚酯8.5\%、尼龙12.5\%等毛利率下，综合毛利约169亿元、毛利率约12\%，与2026Q1 12.88\%方向一致。之后仍需扣除期间费用、利息、税和少数股东；因此收入与毛利检查不能替代归母模型。

\begin{keyinsight}[盈利预测的“所以呢”]
House 77.5亿元不是79.1亿元Street预测的改名，也不是H1年化；它是独立的57.5+20.0亿元桥。\textbf{行动上，H1每偏离中值5亿元先等额调整全年；H2低于18亿元下调基准，H2高于35亿元且现金覆盖资本开支才上调牛情景权重。}
\end{keyinsight}

\section{2027--2028：恢复路径必须有现金和项目证据}

2027E归母70亿元低于2026E，体现异常裂解价差回落；2028E升至75亿元，需要国内加工费、CPL--PA6与部分新项目改善，但不把文莱二期按满负荷纳入。若二期工程与融资快于预期，可在后续模型中增加期权；若煤制乙二醇和二期同时抬高利息及折旧，2028利润反而可能低于当前路径。

\begin{riskbox}[预测失效条件]
若中报实际归母低于55亿元、扣非低于54.7亿元，或H2归母低于18亿元且不是检修／确认时点造成，House基准必须下调；若现金流和短债同时恶化，还应下调PE。\textbf{行动上，不以2027--2028收入增长掩盖2026H2利润坡度。}
\end{riskbox}
